
\chapter{Introduction} 
\noindent Representation theory is an immense topic. Most sources we know
emphasize rigor and depth over conceptual clarity, especially for physics. Such
presentations can conceal the regularity and beauty of the subject just due to
their length. Here, we present the bare minimum, from the physicist's
perspective. We specifically bridge classical physics, where representation
theory is not typically introduced, and quantum physics, where a student
typically first encounters representation theory. We consider this unfortunate,
because representation theory can be useful and clarifying in classical physics
and because quantum physics already has many steep hills to climb without
representation theory. We deem it better to cover representation theory as a
separate topic pertinent to both realms of physics.
